\documentclass[11pt]{article}
\usepackage{mathblatt}
% gesetzt von werkzeuge/setzer.py; Lösungen nach mathe-nachhilfe gemeinsam.md „Lösungen“.
\newcommand{\kennung}{FUR}
\begin{document}
\pfheftstil
\fancyfoot[C]{{\footnotesize\color{mbgrau}\kennung\hfill\thepage}}
\pfheftkopf{Zahlen vergleichen in allen Darstellungen}{Klasse 8 \textperiodcentered{} Lösungen}

\begin{pfloesung}{}
\lz{1b)}{Zehntel}{}{}
\lz{c)}{Tausendstel}{}{}
\lz{d)}{Einer}{}{}
\lz{e)}{Hundertstel}{}{}
\end{pfloesung}
\begin{pfloesung}{}
\lz{2b)}{$2{,}4 > 2{,}38$}{\ (mit Nullen: $2{,}40$ und $2{,}38$)}{}
\lz{c)}{$0{,}09 < 0{,}1$}{\ (mit Nullen: $0{,}09$ und $0{,}10$)}{}
\lz{d)}{$5{,}05 < 5{,}5$}{\ (mit Nullen: $5{,}05$ und $5{,}50$)}{}
\lz{e)}{$1{,}250 = 1{,}25$}{\ (mit Nullen: $1{,}250$ und $1{,}250$)}{}
\end{pfloesung}
\begin{pfloesung}{}
\lz{3.}{$0{,}06 < 0{,}066 < 0{,}6 < 0{,}61$}{$0{,}06$; $0{,}066$; $0{,}6$; $0{,}61$ – mit Nullen: $0{,}060 < 0{,}066 < 0{,}600 < 0{,}610$}{}
\end{pfloesung}
\begin{pfloesung}{}
\lz{4b)}{$-2$}{}{}
\lz{c)}{$-0{,}5$}{}{}
\lz{d)}{$-3$}{}{}
\lz{e)}{$-0{,}1$}{}{}
\end{pfloesung}
\begin{pfloesung}{}
\lz{5b)}{$-1{,}2 < -1{,}02$}{}{}
\lz{c)}{$-0{,}3 < 0{,}1$}{}{}
\lz{d)}{$-2{,}5 < -2{,}45$}{}{}
\lz{e)}{$-0{,}09 > -0{,}1$}{}{}
\end{pfloesung}
\begin{pfloesung}{}
\lz{6b)}{$-0{,}3$}{Abstand $0{,}2$, Hälfte $0{,}1$; $-0{,}4 + 0{,}1 = -0{,}3$}{}
\lz{c)}{$-0{,}8$}{Abstand $0{,}4$, Hälfte $0{,}2$; $-1 + 0{,}2 = -0{,}8$}{}
\lz{d)}{$-0{,}1$}{Abstand $0{,}4$, Hälfte $0{,}2$; $-0{,}3 + 0{,}2 = -0{,}1$}{}
\end{pfloesung}
\begin{pfloesung}{}
\lz{7.}{Kiew, Oslo, Riga, Prag, Wien}{$-4$; $-3{,}5$; $-3{,}05$; $-0{,}8$; $0{,}5$ – also Kiew, Oslo, Riga, Prag, Wien; am kältesten in Kiew.}{}
\end{pfloesung}
\begin{pfloesung}{}
\lz{8b)}{$0{,}75$}{$\frac{3}{4} = \frac{75}{100} = 0{,}75$}{}
\lz{c)}{$0{,}35$}{$\frac{7}{20} = \frac{35}{100} = 0{,}35$}{}
\lz{d)}{$0{,}125$}{$\frac{1}{8} = \frac{125}{1000} = 0{,}125$}{}
\lz{e)}{$0{,}36$}{$\frac{9}{25} = \frac{36}{100} = 0{,}36$}{}
\end{pfloesung}
\begin{pfloesung}{}
\lz{9b)}{$0{,}45$}{$45\,\% = 0{,}45$}{}
\lz{c)}{$0{,}045$}{$4{,}5\,\% = 0{,}045$}{}
\lz{d)}{$1{,}2$}{$120\,\% = 1{,}2$}{}
\lz{e)}{$0{,}005$}{$0{,}5\,\% = 0{,}005$}{}
\end{pfloesung}
\begin{pfloesung}{}
\lz{10b)}{$6\,\% < 0{,}6$}{$0{,}06$ und $0{,}6$, also $6\,\% < 0{,}6$}{}
\lz{c)}{$0{,}35 = \frac{7}{20}$}{$0{,}35$ und $0{,}35$, also $0{,}35 = \frac{7}{20}$}{}
\lz{d)}{$\frac{2}{5} < 45\,\%$}{$0{,}4$ und $0{,}45$, also $\frac{2}{5} < 45\,\%$}{}
\lz{e)}{$\frac{1}{8} < 0{,}13$}{$0{,}125$ und $0{,}13$, also $\frac{1}{8} < 0{,}13$}{}
\end{pfloesung}
\begin{pfloesung}{}
\lz{11b)}{$0{,}04$}{$0{,}2^2 = 0{,}2 \cdot 0{,}2 = 0{,}04$}{}
\lz{c)}{$0{,}25$}{$0{,}5^2 = 0{,}5 \cdot 0{,}5 = 0{,}25$}{}
\lz{d)}{$1{,}21$}{$1{,}1^2 = 1{,}1 \cdot 1{,}1 = 1{,}21$}{}
\lz{e)}{$0{,}01$}{$0{,}1^2 = 0{,}1 \cdot 0{,}1 = 0{,}01$}{}
\end{pfloesung}
\begin{pfloesung}{}
\lz{12b)}{$0{,}3$}{$\sqrt{0{,}09} = 0{,}3$, denn $0{,}3 \cdot 0{,}3 = 0{,}09$}{}
\lz{c)}{$1{,}2$}{$\sqrt{1{,}44} = 1{,}2$, denn $1{,}2 \cdot 1{,}2 = 1{,}44$}{}
\lz{d)}{$0{,}2$}{$\sqrt{0{,}04} = 0{,}2$, denn $0{,}2 \cdot 0{,}2 = 0{,}04$}{}
\lz{e)}{$0{,}7$}{$\sqrt{0{,}49} = 0{,}7$, denn $0{,}7 \cdot 0{,}7 = 0{,}49$}{}
\end{pfloesung}
\begin{pfloesung}{}
\lz{13b)}{$1{,}7 < \sqrt{3} < 1{,}8$}{$1{,}7^2 = 2{,}89 < 3 < 3{,}24 = 1{,}8^2$, also $1{,}7 < \sqrt{3} < 1{,}8$}{}
\lz{c)}{$2{,}2 < \sqrt{5} < 2{,}3$}{$2{,}2^2 = 4{,}84 < 5 < 5{,}29 = 2{,}3^2$, also $2{,}2 < \sqrt{5} < 2{,}3$}{}
\lz{d)}{$0{,}7 < \sqrt{0{,}5} < 0{,}8$}{$0{,}7^2 = 0{,}49 < 0{,}5 < 0{,}64 = 0{,}8^2$, also $0{,}7 < \sqrt{0{,}5} < 0{,}8$}{}
\end{pfloesung}
\begin{pfloesung}{}
\lz{14b)}{$8\,\% = 0{,}08$}{$8\,\% = 0{,}08$; die anderen sind $0{,}9$, $0{,}09$ und $0{,}085$.}{}
\end{pfloesung}
\begin{pfloesung}{}
\lz{15.}{$-\frac{3}{4} < -0{,}7 < 0{,}3 < \sqrt{0{,}16}$}{$-0{,}75$; $-0{,}7$; $0{,}3$; $0{,}4$ – also $-\frac{3}{4} < -0{,}7 < 0{,}3 < \sqrt{0{,}16}$}{}
\end{pfloesung}
\begin{pfloesung}{}
\lz{16.}{$\frac{5}{8} > 0{,}6$}{$\frac{5}{8} = 0{,}625 > 0{,}600$. Falsch: $\sqrt{2} \approx 1{,}41 < 1{,}5$; $0{,}3^2 = 0{,}09 < 0{,}3$; $-0{,}5 > -0{,}6$.}{}
\end{pfloesung}
\begin{pfloesung}{}
\lz{17.}{Laden A; 3 Prozentpunkte mehr als B}{Laden A: ein Viertel $= \frac{1}{4} = 25\,\%$; B: $22\,\%$; C: jeder fünfte Euro $= \frac{1}{5} = 20\,\%$. Am größten bei A, um $25 - 22 = 3$ Prozentpunkte mehr als bei B.}{}
\end{pfloesung}
\begin{pfloesung}{}
\lz{18.}{Ja; Fisch in $3{,}6$\,m Tiefe}{Tiefen als Zahlen: $-4{,}8$ und $-2{,}4$; Abstand $2{,}4$, Hälfte $1{,}2$; $-4{,}8 + 1{,}2 = -3{,}6$. Ja, der Fisch ist in $3{,}6$\,m Tiefe, das ist tiefer als $3{,}5$\,m.}{}
\end{pfloesung}
\end{document}
